\chapter[Sermon 38. Of the Seven Angels Trumpeters, and of the Trumpets: and of the First ii. and iii. Trumpet.]{Of the Seven Angels Trumpeters, and of the Trumpets: and of the First ii. and iii. Trumpet.}
\label{sermon:038}
\markright{Sermon 38. Of the Seven Angels Trumpeters, and of the Trumpets ...}

And the seven angels who had the seven trumpets prepared themselves to blow. The first angel blew: and there was hail and fire mingled with blood, which were cast into the earth; and a third of the trees were burned, and all the green grass withered. And the second angel blew, and as it were a great mountain burning with fire was cast into the sea, and a third of the sea turned to blood; and a third of the creatures which lived in the sea died, and a third of the ships were destroyed. And the third angel blew, and a great star fell from heaven, burning as a torch, and it fell into a third of the rivers, and into the fountains of waters; and the name of the star is called Wormwood, and a third of the waters were turned to Wormwood. And many men died from the waters because they were made bitter.

Our Lord Jesus Christ has kindled on earth a bright and wholesome fire, which the Apostles and those following in their footsteps have continually inflamed. But contrarily, Satan seeks to quench this wholesome fire, and not only to corrupt and deprave this doctrine of salvation, but also to abolish it and overwhelm it with lies. The manner of this is presently described, and even vividly portrayed, for no other end than that the faithful, being warned and fully taught, might be well aware of that pestilent infection. For the scope or end of this book is to preserve the church safe and sound from corruptions, or at least to repair it if it has been corrupted.

Therefore, John saw seven Angels standing in the sight of God. To stand signifies to minister, and encompasses the faith and diligence of ministers. Servants stand before kings, ready to do service and to execute their commandments. We read in the first chapter of Job: “The sons of God came and stood before the Lord, and Satan also appeared among them.” The blessed Angels are called the children or sons of God. They come to do service before God; Satan appears among them, for he too is a minister of God, for the execution of those things that pertain to the wrath and indignation of God against the wicked. All elements are God’s ministers, and finally, all the creatures of God. For he is the Lord of the Sabbath, the God of hosts, who for the salvation and judgment of men uses all his creatures well and rightly, each according to its nature and disposition. For he uses the ministry of Angels like Angels, and so the service of devils as devils. But since the number seven is the number of fullness, containing within itself all times—for there are seven days of creation and rest, there are seven worlds or ages—certainly seven Angels appear before God, for they signify all battles that shall be fought to the end of the world.

For to these seven angels are given seven trumpets, and the angels already had the trumpets, and even prepared themselves to blow the onset. Where chiefly the use of trumpets is to be searched for. The same is most plentifully described of Moses, in the tenth chapter of Numbers. The use of trumpets was diverse, as it is also at this day. First, by the sound of the trumpet, the people of Israel were called together to consult about the common good. Again, at the sound of the trumpet, the senate of princes of the people did assemble. Moreover, they were warned by the trumpet when and who should remove their tents. Furthermore, the trumpets blew to battle, when they joined to fight, as may be seen in the twentieth chapter of Deuteronomy. The people moreover were called together with trumpets on the holy days for public and divine service. “Sound with the trumpet in Zion, call the congregation, gather the people,” says Joel. There was moreover a feast of trumpets, and a Jubilee, having the name of the blowing and sound of trumpets, as appears in the twenty-fifth chapter of Leviticus. Finally, the preaching of the truth was figured by the sound of trumpets, and neither might any other blow it but priests. For it behooves one greatly to whom you commit or deliver the signs public.

Of this varied use of trumpets, none shall agree better with our matter than the warlike. For this world has a shape of war. In it are the camps of good men, and the camps of evil: the tents of Catholics, and the tents of heretics. The chief ruler of these is Satan, and of those Christ: the Captain and Emperor of these is the Devil, of the other the Son of God. And now the angels sound their trumpets and blow the onset: not that the good angels and God himself is the author of heresies and of heretics, whose origin is referred to Satan and sin; but sounding their trumpets they give indeed warning to all men, and signify that most grievous wars shall arise in the world, and even in the church itself. But diverse men are diversely moved and work in war according to their natures. The true Catholics, being warned by the trumpet, take heed to themselves, pray, and finally taking in hand spiritual weapons, prepare themselves unto battle and manfully fight for Christ, and for maintaining and defending the truth. Heretics, sectaries, and men of corrupt minds, according to their malice, taking to them also armor, run forth and fight against Christ and the truth, defend lies, and those who are weak they take, spoil, beat down, and destroy. The good shepherds are the trumpets of God and of Christ: the Devil blows up arch-heretics and beginners of sects.

Regarding the righteous and their struggle, we will hear in chapter 11 and the following chapter. Nevertheless, in every conflict, we must understand that the saints do not sleep, nor are they anywhere idle, but perform their duty everywhere. It was indeed sufficient for the Lord to show us the heretics and sectaries scheming, and to declare how much harm they can inflict, so that we might watch more diligently and beware of all corruption.

\index{Tertullian}\index{Justification}The first angel sounding the first trumpet announces the first conflict. All battles share some likeness and difference. They are alike in that all heresies attack Christ and seek to extinguish or distort the truth of the Gospel. They are diverse in that, at different times, Satan, assailing different doctrines, has spread various heresies throughout the Church. Therefore, as the angel sounds the trumpet, that is to say proclaims war, he warns the saints to be watchful. And as he continues to blow, through God’s permission, according to his just judgment, and by the means and suggestions of Satan, hail and fire mixed with blood were created and sent, or fell, upon the earth. For Paul acknowledges spiritual powers in the heavenly spirits. And Scripture in a certain place aptly depicts wholesome doctrine as a heavy dew and shower that makes the earth fruitful; therefore John rightly compares false and heretical doctrine to hail, for it destroys the fertile places of the earth and utterly ruins the abundant produce. Wherefore, just as elsewhere perverse doctrine is called darnel, leaven, chaff, and so on, it is here called hail. But this hail is tempered and of a wondrous mixture. For it combines fire and blood. These things must be understood allegorically, not literally. Hail is water frozen by cold. And water has been called the wisdom of Scripture; therefore, hail signifies false wisdom. Yet fire is added to it, representing the pretense of Scripture and the inspiration of the Holy Spirit, to which is added blood, the evil affections of humankind—namely, the vices of ambition, wrath, contention, hatred, and similar affections. A hurtful and pestilent doctrine is compounded of these. For when false doctrine governs or corrupts Scripture, and the wicked affections of teachers are joined with it, a pestilent doctrine arises. Such was the doctrine from the beginning of the Nazarenes, or Mimeorites, and of the Hebionites, contending that justification did not come by the alone faith of Christ, but by the law. Our people fought sharply against this pernicious doctrine, namely Paul and the other apostles. From the beginning, some have corrupted the faith with philosophy, others have been blinded by human traditions, and have produced most corrupt opinions. Histories bear witness to this. And Tertullian not without cause called philosophers the patriarchs of heretics. For Paul diligently warned the godly to be wary of philosophy. Those who have not kept themselves from it, and who have valued philosophy and I know not what traditions more highly, have cast great heavy hailstones into the church instead of the heavenly dew and sweet showers.

And have indeed harmed the church greatly. For a third of the trees was burned, and also all the green grass. This number is hinted at in four trumpets, and in six as well. And it seems to signify that a great portion of unstable and wavering people are deceived and abandoned, giving themselves to be destroyed by wicked men; again, the best part of the faithful are to be saved. The Lord himself knows the number exactly. It is enough for us to know these things which he has revealed to us, neither to search curiously any further.

That men are signified by trees appears in the ninth chapter, where it is said, “They had been commanded not to harm the grass of the earth, nor any tree, except men.” After he had said, “Save only those trees which were not marked,” he preferred to say “men,” as if with this key he might unlock the mystery. It is not uncommon in Scripture to represent men by trees, flowers, and grass, as we may gather from Psalm 1, Isaiah 40, and Matthew 12. But that latter point, that all green grass was burned, must be favorably explained. For who could believe that all men were destroyed by those first heresies? We understand therefore that the minds of the faithful were diversely afflicted and tormented by those errors and troubles, but yet, like gold tried in the fire, not to be utterly consumed.

\index{Augustine, Saint}\index{Irenaeus}\index{Epiphanius of Salamis}\index{Daniel (Book of)}\index{Zechariah (Book of)}The second angel sounds the trumpet, signifying that new wars are now brewing; and therefore exhorts that all the godly defend themselves with weapons. And there is cast into the sea not a mountain, but as it were a mountain burning with fire. The sea bears a figure of the world, than which there is nothing more unstable. It is a thing most frequently used by the prophets to call this our world, wherein we live, a sea. By mountains are signified kingdoms, as witnessed by Isaiah in chapter 2, Daniel in chapter 2, and Zechariah in chapter 4, and Christ himself in Matthew 7. By removing of hills or mountains, signifies any difficult thing, and by the opinion of many, an impossible thing. Now therefore springs up an heresy and a doctrine in the church, as it were a burning mountain, which was in deed most furnished, and as it seemed invincible. We read that such was the heresy of the Valentinians, whose sect the holy martyr Irenaeus taught to be divided into many. Such was the fury of the Manicheis and Montanists. They seemed to many to burn with the spirit of God, and to be wholly nothing else but the spirit, and all their oracles to be of the holy ghost. Manichaeus called himself the Apostle of Jesus Christ. The Montanists boasted of a new holy ghost. There was most great plenty of this darnel throughout the universal church. Neither was the success thereof small. For the third part of the sea was made blood. The Apostle signifies the wickedness of sects. For how vile and impudent were the heretics called Gnostics, the Valentinians, and Manicheis, as testified by Irenaeus, Augustine, and Epiphanius. And a great part of the creatures in the sea perished. And he speaks of those who have souls, not of fishes in deed, but men. Many ships moreover were lost, to wit, mariners and land men, being corrupted with these heresies.

Those heresies arose from the writings of the authors I mentioned, but they are not yet completely gone; corrupt men persistently associate with one another, reviving the old errors. Consequently, a bitter conflict remains in the church to this day, and we are continually warned to be on guard against these corruptions.

\index{Jerome, Saint}The third angel blows his trumpet, proclaiming new wars; and behold, a great star fell from heaven, burning like a torch, and infected the third part of the rivers and fountains of waters. That star is called wormwood. I told you in the first chapter that stars are called preachers, bishops, and notable men in the church. Therefore, it signifies that some notable man should fall away from the true faith into heresy, with which he should infect a great part of the world, corrupting the Scriptures and the sound doctrine of faith. And these things seem to be fulfilled in Paulus Samosatenus and Arius. This torch burned horribly and inflamed the whole world without recovery. That pestilence denied the deity of Christ and made the whole Gospel most bitter to us. For if Christ is not very God, how is he a Savior, King, bishop, intercessor, mediator, and salvation of the faithful? He quenched the light that denied the deity of Christ. Therefore, he is called by the name of wormwood. The prophet Jerome used the selfsame allegory, or metaphor, or allusion, in the ninth and twenty-third chapters. And Amos in the sixth chapter, where he says that the judges have turned judgment into wormwood.

\index{Eusebius of Caesarea}\index{Zurich}The Scripture and doctrine, beautifully depicted by rivers and fountains, was the occasion of death for many, having been corrupted by the Arians. The Scripture and doctrine of the Gospel is of itself mortal to no one, but rather lively to all; yet corruption makes it deadly. Poison put in wine makes the wine deadly: the wine itself kills no one, but rather glads and rejoices all people. Read the ecclesiastical histories of Eusebius, Theodoret, Sozomenus, Socrates, and others, and you shall perceive how aptly John has written all these things, and how rightly they are all fulfilled. No small part of that bitterness has flowed unto our time, while that old error is often renewed by the instigation of the devil. For what that unclean beast, Michael Servetus, a Spaniard, vomited against the Son of God, for his impenitent wickedness and continual blasphemy, the world knows; he was burned at Zurich. We must therefore pray to the Lord that in such dangerous conflicts he would keep us safe and sound. Amen.
